\begin{ouverture}[college2]{9}{Décrire et mesurer un mouvement}{Partie 2 · Mouvements et interactions}{Un passager assis est-il immobile ?}
\objectif{Choisir un objet et un référentiel.}
\objectif{Décrire trajectoire, direction, sens et évolution de la vitesse.}
\objectif{Calculer distance, durée et vitesse avec des unités cohérentes.}
\prerequis{Durées, distances et proportionnalité.}
\end{ouverture}
\pcPage{Le cours}{Décrire et mesurer un mouvement}
\phantomsection\label{ess-9}
\essSection{1}{Le mouvement dépend du référentiel}
Pour décrire un mouvement, on choisit l’objet étudié et un référentiel, c’est-à-dire un objet de référence auquel on associe des repères d’espace et de temps. Un passager assis est immobile par rapport à son siège, mais en mouvement par rapport au quai lorsque le train roule. Ces descriptions sont compatibles parce qu’elles utilisent des références différentes.\par
\pcFigure{figures/c09-chrono.pdf}{Des positions relevées à intervalles de temps égaux : l’écartement renseigne sur l’évolution de la valeur de la vitesse.}
\essSection{2}{Décrire la trajectoire et la vitesse}
La trajectoire est l’ensemble des positions successives du point étudié dans le référentiel choisi. Elle est rectiligne, circulaire ou plus généralement curviligne. La vitesse possède une direction, un sens et une valeur. Un mouvement circulaire à valeur de vitesse constante change néanmoins de direction. « Uniforme » signifie que la valeur de la vitesse reste constante ; cela ne signifie pas forcément rectiligne.\par
\clearpage
\essSection{3}{Relier distance, durée et vitesse}
La vitesse moyenne sur un trajet est $v_{\mathrm{moy}}=d/\Delta t$. Pour un mouvement uniforme, on utilise $d=v\Delta t$ et $\Delta t=d/v$. Si $d$ est en m et $\Delta t$ en s, $v$ est en m/s. Pour passer de m/s à km/h, multiplier par $3{,}6$ ; pour le passage inverse, diviser par $3{,}6$. Une vitesse moyenne ne décrit pas tous les instants du trajet.\par
\pcFigure{figures/c09-graphique.pdf}{Sur ce graphique de position rectiligne, une pente plus forte correspond à une valeur de vitesse plus élevée.}
\essSection{4}{Mesurer sans inventer une précision}
Une chronophotographie ou une vidéo fournit des positions et des instants. Pour exploiter l’image, il faut une échelle de distance, une durée entre images et un point suivi cohérent. La perspective peut déformer les longueurs. Une mesure avec un chronomètre manuel comporte une variabilité due au déclenchement. Répéter et comparer aide à estimer la qualité du résultat.\par
\begin{pcEncadre}{method}{Un exemple expliqué}Une cycliste parcourt $180\unite m$ en $30\unite s$. Sa vitesse moyenne vaut $180/30=6{,}0\unite{m/s}$, soit $6{,}0\times3{,}6=21{,}6\unite{km/h}$. Si elle conserve cette vitesse pendant $50\unite s$, elle parcourt $d=6{,}0\times50=300\unite m$. La dernière conclusion suppose bien un mouvement uniforme sur la durée annoncée.\end{pcEncadre}
\pcPage{Expérimenter et retenir}{Décrire et mesurer un mouvement}
\phantomsection\label{activite-9}
\begin{pcEncadre}{method}{Investigation en classe}\textbf{Question.} Un petit chariot avance-t-il uniformément ? Filmer un déplacement rectiligne sur une table, avec un repère de longueur dans le plan du mouvement. Choisir le centre d’une même roue ou un point repéré du chariot. Relever les positions toutes les mêmes durées, puis tracer position en fonction du temps. Comparer les distances successives et discuter les erreurs de lecture. \textbf{À rendre :} référentiel, tableau, graphique et conclusion. Aucune mesure ne se fait sur une route ouverte.\end{pcEncadre}
\pcSousTitre{Ce que je dois savoir expliquer}\begin{questions}
\item Un mouvement est décrit dans un référentiel annoncé.
\item Trajectoire et évolution de la vitesse sont deux descriptions complémentaires.
\item À temps égaux, des déplacements égaux indiquent une valeur de vitesse constante.
\item La vitesse moyenne vaut distance totale divisée par durée totale.
\item Les unités doivent être converties avant le calcul.
\end{questions}
\begin{pcEncadre}{attention}{Une erreur à éviter}Faire la moyenne de deux vitesses ne donne pas toujours la vitesse moyenne du trajet : il faut utiliser la distance totale et la durée totale.\end{pcEncadre}
\pcRepere{cor-activite-9}{Pistes pour l’investigation}\pcGrille
\clearpage\markright{Je démarre}
\rubrique{Je démarre}
\par\noindent\begin{minipage}{\linewidth}
\begin{exercice}[1]{Le référentiel}
\begin{questions}
\item Un livre posé dans un train est-il immobile par rapport au train ?
\item Peut-il être en mouvement par rapport au sol ?
\end{questions}
\end{exercice}
\end{minipage}
\par\noindent\begin{minipage}{\linewidth}
\begin{exercice}[1]{La trajectoire}
\begin{questions}
\item Un point décrit un cercle autour d’un axe. Comment qualifier la trajectoire ?
\item Un mouvement en ligne droite est-il circulaire ?
\end{questions}
\end{exercice}
\end{minipage}
\par\noindent\begin{minipage}{\linewidth}
\begin{exercice}[1]{Temps égaux}
\pcFigureEx{figures/ex-chronophoto.pdf}{Les prises de vue sont séparées par une même durée.}
\begin{questions}
\item Sur une chronophotographie, les écarts augmentent pour des durées égales. Que peut-on dire de la valeur de la vitesse ?
\item Quelle donnée faut-il vérifier ?
\end{questions}
\end{exercice}
\end{minipage}
\par\noindent\begin{minipage}{\linewidth}
\begin{exercice}[1]{Une vitesse directe}
\begin{questions}
\item Un mobile parcourt $40\unite m$ en $8\unite s$. Calculer sa vitesse moyenne.
\item Donner l’unité.
\end{questions}
\end{exercice}
\end{minipage}
\par\noindent\begin{minipage}{\linewidth}
\begin{exercice}[1]{Convertir}
\begin{questions}
\item Convertir $10\unite{m/s}$ en km/h.
\item Convertir $72\unite{km/h}$ en m/s.
\end{questions}
\end{exercice}
\end{minipage}
\clearpage\markright{Je m’entraîne}
\rubrique{Je m’entraîne}
\par\noindent\begin{minipage}{\linewidth}
\begin{exercice}[2]{Calculer une distance}
\begin{questions}
\item Un chariot roule à $0{,}40\unite{m/s}$ pendant $15\unite s$ à vitesse constante. Quelle distance parcourt-il ?
\item Et pendant une durée double ?
\end{questions}
\end{exercice}
\end{minipage}
\par\noindent\begin{minipage}{\linewidth}
\begin{exercice}[2]{Calculer une durée}
\begin{questions}
\item Une marcheuse parcourt $1{,}8\unite{km}$ à $6\unite{km/h}$ de moyenne. Calculer la durée en heures.
\item Convertir en minutes.
\end{questions}
\end{exercice}
\end{minipage}
\par\noindent\begin{minipage}{\linewidth}
\begin{exercice}[2]{Lire des positions}
\pcFigureEx{figures/ex-position.pdf}{}
\begin{questions}
\item Aux instants 0, 2, 4 et 6 s, les positions sont 0, 3, 6 et 9 m. Décrire le mouvement rectiligne.
\item Calculer la vitesse.
\end{questions}
\end{exercice}
\end{minipage}
\par\noindent\begin{minipage}{\linewidth}
\begin{exercice}[2]{Uniforme sur un cercle}
\begin{questions}
\item Un siège de manège décrit un cercle à valeur de vitesse constante. Sa direction de vitesse reste-t-elle fixe ?
\item Comment qualifier le mouvement ?
\end{questions}
\end{exercice}
\end{minipage}
\par\noindent\begin{minipage}{\linewidth}
\begin{exercice}[2]{Une vidéo étalonnée}
\begin{questions}
\item Une règle de $1\unite m$ occupe $200$ pixels ; un déplacement occupe $80$ pixels. Quelle distance réelle ?
\item Ce déplacement dure $0{,}20\unite s$. Calculer la vitesse moyenne.
\end{questions}
\end{exercice}
\end{minipage}
\clearpage\markright{Je raisonne}
\rubrique{Je raisonne}
\par\noindent\begin{minipage}{\linewidth}
\begin{exercice}[3]{Deux parties d’un trajet}
\begin{questions}
\item Un trajet de $60\unite{km}$ dure $1\unite h$ ; le retour de $60\unite{km}$ dure $2\unite h$. Calculer la vitesse moyenne de l’aller-retour.
\item Pourquoi $45\unite{km/h}$ serait-il faux ?
\end{questions}
\end{exercice}
\end{minipage}
\par\noindent\begin{minipage}{\linewidth}
\begin{exercice}[3]{Un arrêt compte-t-il ?}
\begin{questions}
\item Un bus parcourt $12\unite{km}$ en $20\unite{min}$, puis s’arrête $10\unite{min}$. Calculer sa moyenne sur les 30 minutes.
\item La vitesse pendant l’arrêt vaut-elle cette moyenne ?
\end{questions}
\end{exercice}
\end{minipage}
\par\noindent\begin{minipage}{\linewidth}
\begin{exercice}[3]{Une conclusion trop forte}
\begin{questions}
\item Une voiture parcourt 100 m en 5 s. Peut-on prouver qu’elle était à 20 m/s à chaque instant ?
\item Que permet réellement le calcul ?
\end{questions}
\end{exercice}
\end{minipage}
\par\noindent\begin{minipage}{\linewidth}
\begin{exercice}[3]{Un graphique ambigu}
\pcFigureEx{figures/ex-position-retour.pdf}{Exemple de position décroissant linéairement.}
\begin{questions}
\item Une courbe représente une position qui diminue linéairement au cours du temps. Le mobile ralentit-il forcément ?
\item Quelle distinction faut-il faire ?
\end{questions}
\end{exercice}
\end{minipage}
\par\noindent\begin{minipage}{\linewidth}
\begin{exercice}[3]{Comparer des mesures}
\begin{questions}
\item Trois temps pour $10\unite m$ valent $4{,}8$, $5{,}0$ et $5{,}2\unite s$. Calculer le temps moyen.
\item Donner une vitesse représentative et discuter sa précision.
\end{questions}
\end{exercice}
\end{minipage}
\pcPage{Défi du chapitre}{Décrire et mesurer un mouvement}
\phantomsection\label{defi-ex-9}
\pcDefiTitre{9}{Choisir le meilleur trajet}
\begin{questions}
\item A mesure $1{,}2\unite{km}$ à $4\unite{km/h}$ de moyenne. B mesure $1{,}5\unite{km}$ à $6\unite{km/h}$, avec $3\unite{min}$ d’arrêt supplémentaire. Calculer les durées totales.
\item Quel trajet est le plus rapide ?
\item Quelle hypothèse pourrait changer la conclusion ?
\end{questions}
\pcGrille
